\section{A checker that survives lowering}
\label{part:lowering}

Everything up to here analyses \emph{one} artifact. This part asks what happens
when that artifact is compiled further, and derives the architecture the
measurements force. The short version: single-point analysis is not salvageable,
and the fix is to stop pushing policy downward.

\subsection{What the measurements found}

Twenty semantics-preserving transformations were reproduced against LLVM 17.0.6.
They sort into three gaps, and it is worth naming them separately because they
fail in different directions.

\paragraph{The level gap.} The transformation happens \emph{below} LLVM IR, so an
IR-level verdict describes a program that does not run. The sharpest instance:
\code{x86-cmov-converter} turns a conditional move back into a branch. Composed
with mid-end speculation, which turns a branch into a select, the pipeline
produces the worst possible sequence --- \textbf{leaky source, clean \code{-O2}
IR, leaky binary}. The same IR yields opposite verdicts on x86-64,
aarch64-generic, aarch64-neoverse-n2, and riscv64. Related: \code{StackColoring}
merges two lifetime-bracketed allocations onto one frame offset, with
runtime-confirmed residue.

\paragraph{The annotation gap.} The policy does not survive. \code{combineMetadata}
discards custom metadata when merging; every signature-changing pass drops
parameter attributes; \code{mlir-translate} drops MLIR attributes entirely, with
no diagnostic. Combined with the convention that an unlabelled value is public,
each of these is a false negative rather than a refusal.

\paragraph{The counting gap.} Occurrence cardinality moves in \emph{both}
directions under plain \code{-O2} --- sinking merges two sites into one, jump
threading splits one into two. So no exact-count policy is an artifact
invariant.

\begin{remark}[The finding underneath all twenty]
\label{rem:notnormative}
The artifact the checker reads is not the artifact that runs, and nothing in the
system says which one is normative. Every other problem in this part is a
consequence.
\end{remark}

\subsection{Why single-point analysis cannot be repaired}

There are only two places to put a single analysis, and the measurements close
both.

\textbf{Analyse early}, where the policy is meaningful and the structure is
intact. Then the level gap applies in full: the binary can leak through a channel
that does not exist at that level, and the laundering sequence above produces a
clean verdict on a leaky program.

\textbf{Analyse late}, close to the machine, where the observations are real.
Then the annotation gap applies: the policy is gone. Worse, it is gone
\emph{silently}, so the analysis does not refuse --- it sees unlabelled code and,
under the usual convention, calls it public.

The obvious repair is to carry the policy downward so that late analysis has
something to read. \cref{sec:carriers} shows every known carrier has a silent-loss
path. That is what makes this an architectural problem rather than an engineering
one.

\subsection{Every carrier leaks}
\label{sec:carriers}

\begin{center}
\begin{tabular}{@{}p{0.30\linewidth} p{0.60\linewidth}@{}}
\toprule
Carrier & Measured failure \\
\midrule
MLIR \code{sps.*} attributes
  & \code{mlir-translate --mlir-to-llvmir} drops all of them; zero survive, no
    diagnostic \\
LLVM metadata
  & \code{combineMetadata} discards unknown kinds when merging two sites \\
Parameter attributes
  & dropped by every pass that rewrites a signature \\
Opaque marker call
  & must have no visible definition to survive, but then it does not link;
    once the definition is visible it is inlined away \\
Debug info
  & explicitly non-semantic; passes may drop or coarsen it \\
\bottomrule
\end{tabular}
\end{center}

Note what these have in common. They are all attempts to make a \emph{policy
fact} travel alongside the code, through a pipeline whose passes have no
obligation to preserve anything except semantics. The pipeline is not
misbehaving; it was never told these facts mattered.

\subsection{The inversion}
\label{sec:inversion}

If policy cannot travel down, send observations up.

\begin{center}
\begin{tikzpicture}[
  lvl/.style={draw, rounded corners=2pt, minimum height=7mm, minimum width=26mm,
              font=\small, inner sep=3pt},
  lbl/.style={font=\scriptsize},
  ar/.style={-{Latex[length=1.8mm]}}
]
  \node[lvl, fill=blue!8] (src) at (0,2.1)  {source / MLIR};
  \node[lvl]              (ir)  at (0,1.05) {LLVM IR};
  \node[lvl]              (mir) at (0,0)    {machine IR};
  \node[lvl]              (bin) at (0,-1.05){binary};

  \draw[ar] (src) -- (ir);   \draw[ar] (ir) -- (mir);   \draw[ar] (mir) -- (bin);

  \node[lbl, anchor=west, text width=4.3cm] at (1.9,2.1)
    {\textbf{policy lives here} --- authored, meaningful, stable};
  \node[lbl, anchor=west, text width=4.3cm] at (1.9,-0.5)
    {\textbf{observations live here} --- branches, addresses, frame slots, event order and count};

  \draw[ar, dashed] (3.4,-0.15) .. controls (5.2,0.9) .. (3.4,1.95);
  \node[lbl, anchor=west] at (4.6,0.95) {provenance};
\end{tikzpicture}
\captionof{figure}{Policy stays at the level where it is authored. Observations
are extracted at the level where they are real. A provenance map relates the
second to the first, and the relational check runs once, at the top, over
observations sourced from the bottom.}
\label{fig:inversion}
\end{center}

The three gaps close for structural reasons rather than by better engineering:

\begin{itemize}[nosep,leftmargin=1.4em]
\item \textbf{Level gap} --- observations are gathered where they are real, so a
      branch introduced by the backend is simply present in the observation set.
\item \textbf{Annotation gap} --- labels never travel, so nothing can drop them.
\item \textbf{Counting gap} --- cardinality is \emph{measured} at the bottom
      rather than asserted at the top, so a merged or duplicated site is
      observed, not assumed away.
\end{itemize}

\subsection{The quantifier error, and universal discharge}
\label{sec:quantifier}

The architecture as first stated has a defect worth dwelling on, because it is a
\emph{logic} error rather than an engineering gap, and because the first draft of
this section contained it.

That draft said: \emph{for each policy site, check the predicate over the
observations attributed to it}. That quantifies existentially over sites.
Soundness requires universal quantification over \emph{observations}.

The difference is not pedantic. A per-site predicate cannot distinguish
\textbf{``no observation here''} from \textbf{``compliant here''}. So every
mechanism that moves an observation off a named site --- outlining,
misattribution, a build without debug info --- leaves a site with nothing
attached, which the per-site rule reads as compliance. It \emph{fails open},
silently.

Reproduced with extractor, provenance map and artifact all held fixed: one
aarch64 build at \code{-Oz} yields \Verified{} under the per-site rule and a
refusal under the corrected one.

\paragraph{The correction.} \textbf{Every observation of a policy-relevant kind
in the artifact must be discharged by exactly one policy site.} Undischarged
observations are the failure condition; per-site predicates run only \emph{after}
that obligation is met.

The cost was measured, not estimated: across 42 fixtures $\times$ 2 targets
$\times$ 2 optimization levels, 167 of 168 builds admissible --- a refusal
surface near 1\%. It is affordable because it quantifies over policy-relevant
observation \emph{kinds}, not over instructions.

\begin{remark}[The same mistake, one level down]
The original problem was that policy could not survive travelling downward. The
first fix moved policy out of the way and rested the decision on the \emph{site
namespace} instead --- but a site namespace is another thing the compiler is free
to rewrite, so that repeats the original error against a different object.
Provenance is demoted accordingly: it may localize a finding, never decide one.
\end{remark}

\subsection{What each stage owes}

\begin{center}
\begin{tabular}{@{}l p{0.62\linewidth}@{}}
\toprule
Stage & Obligation \\
\midrule
Policy (top)
  & declare secrets, observer, releases, placement. Never travels down. It may be
    \emph{composed} across translation units: that is composition at the top and
    crosses no level \\
Extraction (each level)
  & emit the observation set for that level. This is what earns the
    architecture: observations existing at no higher level, such as an outlined
    function's secret-indexed loads, are seen only here \\
Provenance
  & a \textbf{localization hint}, not a decision mechanism. Three-valued:
    present, absent, or present-but-unattested. The third state is load-bearing
    --- confidently wrong provenance must degrade a diagnostic, never ship a
    leaky artifact marked verified \\
Decision (top)
  & \textbf{universal discharge}: every policy-relevant observation accounted for
    by exactly one site. Per-site predicates run afterwards \\
Differential
  & the actual soundness argument, not a supplement. Runs over \emph{whole
    images}, never over policy-named sites, under one composed program-level
    policy \\
Refusal
  & absent provenance, unattested provenance, absent callee model, or any
    undischarged observation \\
Identity
  & binds \emph{both} ends: input closure, link unit, target triple, subtarget
    features, codegen flags, compiler revision \\
\bottomrule
\end{tabular}
\end{center}

\begin{remark}[When this is worth the complexity, and when it is not]
\label{rem:simpler}
For a \emph{single pinned target tuple} a simpler architecture is equally sound
--- analyse the final artifact only, and re-supply policy at that level through a
side file keyed on symbol and offset rather than carried in the IR. At one site,
stack-slot merging, it is strictly \emph{more} sound.

The staged architecture earns its complexity only when more than one target tuple
ships. That is a checkable condition, and it should be checked before paying the
cost rather than assumed.
\end{remark}

\begin{remark}[The role of markers changes, and shrinks]
Under \cref{sec:inversion} a declassification marker no longer has to survive to
the binary. It only has to survive to the level where the \emph{policy} is
evaluated. That is a far weaker requirement than the one the attack findings
demolished, and it is why the inversion is worth the engineering: it converts a
problem with no known sound solution into one with a workable one.
\end{remark}

\subsection{The declassification marker, after the attacks}
\label{sec:marker}

Markers still have a job: recording that a High-to-Low transition was
\emph{authorized}. The obvious realization --- an opaque function the optimizer
cannot see through --- was attacked across five dimensions and does not survive.
The failure is structural, not a tuning problem.

\paragraph{Why the function-shaped marker fails.} A declaration must resolve at
link time. The moment its definition is co-visible --- LTO, ThinLTO, a header,
the same translation unit --- attribute inference deduces \code{returned} from
the identity body, the optimizer replaces uses of the marker's result with its
argument, and the release edge is severed \emph{while the marker sits there
decoratively}. A presence-based checker sees a marker and a clean path; the raw
secret reaches the sink. Verified end-to-end through the real \code{clang} driver
under full LTO.

That is the whole family: no definition means no link, and a definition means no
protection.

\paragraph{What works instead: tied-register inline asm.} There is no symbol, so
there is nothing to define, nothing to link, and nothing for LTO to import.

\begin{center}
\begin{minipage}{0.92\linewidth}
\small
\begin{verbatim}
%h = call T asm sideeffect "// sps.declassify id=$2",
                           "=r,0,i,~{memory}" (T %raw, i32 <ID>)
\end{verbatim}
\end{minipage}
\end{center}

Each operand earns its place, and each was measured:

\begin{center}
\begin{tabular}{@{}l p{0.68\linewidth}@{}}
\toprule
Operand & What it buys \\
\midrule
\code{sideeffect}
  & prevents deletion when the result is unused, and prevents machine-level CSE
    merging two markers \\
\code{=r} output
  & puts the marker \emph{on} the dataflow path. With no body, \code{returned}
    cannot be inferred --- the exact failure of the function-shaped form \\
\code{0} tie
  & ties input to output, so the machine cost is zero. Measured 9 instructions
    marked versus 9 unmarked, no frame, no spills, tail call preserved \\
\code{i} on the id
  & the id must be an immediate. A non-constant id is a \textbf{hard
    instruction-selection error}, so identity recoverability is fail-closed by
    construction rather than by convention \\
\code{\textasciitilde\{memory\}}
  & no IR-level effect, but blocks post-register-allocation scheduling from
    reordering memory operations across the marker. Measured cost zero \\
\bottomrule
\end{tabular}
\end{center}

Verified identical under \code{default<O1/O2/O3/Os/Oz>}, \code{lto<O2>},
\code{lto<O3>} and \code{thinlto<O2>}: marker present, id literal, sink consuming
the marker's result --- while the declassifier implementation is fully inlined
and deleted.

\paragraph{The freeze point this forces.} Post-LTO, post-optimization IR,
immediately before instruction selection. The asm comment is dropped by the
assembler --- measured one occurrence in the \code{.s}, zero in the \code{.o},
zero in the linked binary --- so anything later cannot see the marker at all.
Where a shipped-binary attestation is genuinely required, the comment body can be
replaced by a section note emitting a \code{(pc, release\_id)} record with zero
text instructions.

\paragraph{Two problems that do not go away.} Stated as residual risk rather than
buried:

\begin{itemize}[nosep,leftmargin=1.4em]
\item \textbf{\code{returned} is still expressible.} Inline asm does not
      structurally forbid it --- the IR verifier accepts \code{returned} on an
      asm argument, and the bypass reproduces. It must be an explicit checker
      rule, because there is no IR-level prohibition.
\item \textbf{A static marker count is not a release count.} Unrolling and
      inlining both multiply it. Occurrence cardinality has to be restated
      dynamically, or checked at an earlier freeze point --- which is the
      counting gap of \cref{rem:notnormative} reappearing, unsolved.
\end{itemize}

\begin{remark}[The general lesson]
The function-shaped marker failed because it asked the \emph{linker} and the
\emph{optimizer} to respect a policy fact neither was told about. The inline-asm
marker works because it asks for nothing: it is opaque by construction, not by
permission. Every durable mechanism in this part has that shape --- it survives
because there is no decision point at which something could decide otherwise.
\end{remark}

\subsection{The differential check}

Even with the inversion, run the analysis at each level independently and compare
verdicts. Disagreement between levels is itself the finding:

\begin{center}
\begin{tabular}{@{}l l p{0.44\linewidth}@{}}
\toprule
Upper & Lower & Meaning \\
\midrule
\Verified{} & \Unsafe{}   & \textbf{laundering} --- the lowering introduced the leak.
                            No fixture in the corpus covers this direction \\
\Unsafe{}   & \Verified{} & evidence destroyed; the leak was optimised out of view \\
\Verified{} & \Unknown{}  & acceptable: monotone degradation \\
\Unknown{}  & \Verified{} & suspicious --- what fact appeared? \\
\bottomrule
\end{tabular}
\end{center}

This is cheap. It needs no new theory, it reuses the existing corpus, and it is
the mechanism that would have caught the inlined-declassifier case automatically
rather than by inspection.

\subsection{What this costs, and what it does not buy}

The honest accounting.

\textbf{It costs backend work.} Extraction at machine-IR level is not an IR pass;
it means reading or instrumenting the target layer, per target. The measurements
already show four target tuples disagreeing on one input, so this is per-tuple
work.

\textbf{Provenance is best-effort.} Debug-info \code{inlinedAt} chains do relate
inlined instructions back to their original function --- verified --- but debug
info is explicitly non-semantic and may be dropped. So the provenance map is a
source of \Unknown{}, not a guarantee. That is acceptable precisely because
\Unknown{} is a first-class outcome; it would be unacceptable in a three-valued
design.

\textbf{It is not a proof.} This architecture makes the checker's answer be about
the artifact that runs. It does not establish that the answer is right --- that
is still \cref{thm:core} and the assurance ladder of
Part~\ref{part:correctness}.

\textbf{It does not close \Lv{4}.} Real latency, hardware isolation, and
speculation remain outside any observation set you can extract statically.

\subsection{One capture point, several extraction points}
\label{sec:capture}

Two results in this document appear to disagree about where the artifact is
frozen, and the disagreement is terminological. Resolving it fixes the most
important architectural parameter, so it is worth doing explicitly.

The marker analysis concludes the freeze point is \textbf{post-LTO, immediately
before instruction selection}, because the marker is not visible below that ---
measured, one occurrence in the \code{.s}, zero in the \code{.o}, zero in the
linked binary. The inversion concludes observations must come from \textbf{the
machine level}, because that is where \code{je} rather than \code{csel} is
decided. Both are right. They are answers to different questions that were being
given one name.

Split the concept:

\begin{center}
\begin{tabular}{@{}l p{0.62\linewidth}@{}}
\toprule
\textbf{Capture point} & where the \emph{annotated} artifact is snapshotted and
hashed. Exactly one. Must be where policy is still readable: post-LTO, before
instruction selection. Markers must survive to here and no further. \\[3pt]
\textbf{Extraction points} & where \emph{observations} are gathered. Several, at
and below the capture point, including machine IR. Carry no policy at all. \\[3pt]
\textbf{Artifact identity} & binds the capture point to every extraction point,
so the report describes one compilation rather than a family. \\
\bottomrule
\end{tabular}
\end{center}

The requirement on markers weakens accordingly, and this is the practical
payoff: a marker must survive to the capture point, not to the binary. That is a
far easier obligation, and it is why \cref{sec:marker}'s inline-asm form is
sufficient despite the assembler discarding it.

\begin{remark}[What the identity must bind]
Binding the top-level module alone is not enough --- the same module produced
\code{je} on x86-64 and \code{csel} on aarch64 from identical IR. So the identity
is a \emph{tuple}: module hash, ordered pipeline digest, target triple, target
CPU and features, and optimization level. A verdict that does not name that tuple
is not a claim about anything runnable. This is the same shape of requirement as
per-coalition rows, and it lands on the same record.
\end{remark}

\subsection{Open problems}

\paragraph{Irreducible: event multiplicity.} This one is not an open problem to
be solved later; it is a limit of the architecture, and naming it is the only
honest handling.

Provenance can be \emph{perfect} and the count still wrong. Measured: at
\code{-O2 -g}, two \code{stp} store-pair instructions carrying entirely correct
source attribution perform sixteen release events --- vectorization at four to
one, composed with store-pairing at two to one. Nothing is misattributed. The
provenance map is exactly right. The count is off by a factor of eight.

Worse, the defect is structural rather than incidental. The decision quantifies
over the \emph{union} of observation sets, and set union is idempotent, so the
formulation cannot count at all. A policy constraining release cardinality is
therefore not checkable in this architecture as stated. It needs either a
multiset, which changes the decision procedure, or a dynamic statement, which
changes what is being claimed.

Until one of those is chosen, a release-count policy must be reported
\Unknown{}, not approximated.

\medskip
Two further problems remain open, stated rather than hidden.

\textbf{Non-monotone cardinality.} Because occurrence counts move both ways, a
policy that constrains release count cannot be checked as an artifact invariant
at a single level. It has to be a statement about the observation set, which
means the extraction has to be trusted to be complete.

\textbf{Multi-target verdicts.} One artifact and four targets gave four different
answers. Either a verdict names its target tuple, or the artifact must be
verified for every tuple it will be built for --- and the report needs a row per
tuple, which is the same record-shape problem as coalitions.
